\section{Memory Technologies}
\label{sec:memory_technologies}

There are two types of external memory which are relevant to this work: DDR2 SDRAM, which provides large-capacity volatige storage, and Flash memory, which provides non-volatile storage for persistent data. Their characteristics differ significantly, and understanding these differences is important for understanding architectural design choices described in later chapters.

\subsection{DDR2 SDRAM}
\label{sec:ddr2}

Double Data Rate 2 (DDR2) Synchronous Dynamic Random-Access Memory (SDRAM) is a memory standard which can transfer data on both the rising and falling edges of a clock signal, doubliing the data rate compared to single data rate memories \textcolor{red}{source for DDR2}. It offers large and inexpensive storage capacity, however due to its volatile nature the data stored does not persist when powering off.

\textcolor{red}{fig showing SDR vs DDR?}

The access latency in DDR2 memory is variable, and depends on wether the requested address falls within the currently open row in a memory bank. A row hit is accessed at relatively low latency, while a row miss requires a precharge and activate sequence before data can be accessed, which results in significantly higher latency \textcolor{red}{source}. This makes DDR2 well suited for sequential access patterns, but less efficient for random access.

\subsection{Flash Memory}
\label{sec:flash}

Flash memory provides non-volatile storage, meaning it retains data contents when powered off \textcolor{red}{source}. This makes it suitable for storing data which should survive a system reset, like for example trained model parameters.  Read operations are relatively fast and byte-addressable, however read throughput is still far slower than DDR2. Write and erase operations are significantly slower than reads, and erasure must be performed on entire sectors at a time. In addition to this the Flash has only a limited number of lifetime write and erase operations, meaning high frequency changes in stored data is not recommended, and the Flash is better suited as read-only memory (ROM).

The Quad SPI (QSPI) interface is a common way to connect Flash devices to FPGAs and microcontrollers, which uses four parallel data lines to  provide higher read throughput than standard SPI. Most devices however also support the standard SPI interface, and double SPI interface. \textcolor{red}{sources sources and more sources}
